% Generated from locale/tr/content/computability/computability-theory/normal-form.tex; do not hand-edit.


\olfileid[tr]{cmp}{thy}{nfm}
\olsection{Normal Biçim Teoremi}

Hesaplanabilir fonksiyonların tanımlarını betimleyebildiğimizi ve bu
fonksiyonlardan birinin belirli bir girdideki hesabının tam kaydına ilişkin
varsayılan bir betimin doğru olup olmadığını sınayabildiğimizi düşünelim.
O zaman hesaplama modelinden bağımsız olarak, hesaplanabilir herhangi bir
$f$ fonksiyonunun herhangi bir $x$ girdisindeki değerini şu yolla
belirleyebilmemiz beklenir:
\begin{enumerate}
\item Hesap kayıtlarının bütün olası betimleri arasında arama yapmak.
\item Belirli bir kaydın $f(x)$ hesabının kaydı olup olmadığını sınamak.
\item Doğru kayıt bulunmuşsa $f(x)$ değerini bu kayıttan çıkarmak.
\end{enumerate}
Bunun gerçekten doğru olması Kleene normal biçim teoreminin içeriğidir.

\begin{thm}[Kleene Normal Biçim Teoremi]
\ollabel{thm:normal-form}
Şu özelliği taşıyan ilkel özyinelemeli bir $T(e,x,s)$ bağıntısı ve ilkel
özyinelemeli bir $U(s)$ fonksiyonu vardır: $f$ herhangi bir kısmi
hesaplanabilir fonksiyonsa bazı $e$ için her $x$'te
\[
  f(x)\simeq U(\umin{s}{T(e,x,s)})
\]
olur.
\end{thm}

\begin{proof}[İspat Taslağı]
Her hesaplama modeli için hesaplanabilir $f$ fonksiyonunun bir betimi titiz
biçimde tanımlanıp böyle bir betim $e$ doğal sayısıyla kodlanabilir. Ayrıca
indisi $e$ olan fonksiyonun $x$ girdisindeki hesap sürecini kaydeden bir
``hesap dizisi'' kavramı da titiz biçimde tanımlanabilir. Böyle bir hesap
dizisi de $s$ sayısıyla kodlanabilir. Kodlama şu iki nesne hesaplanabilir
olacak biçimde yapılabilir:
\begin{enumerate}
\item Bir $s$ sayısı indisi $e$ olan fonksiyonun $x$ girdisindeki hesap
  dizisini kodladığında ve ancak o zaman geçerli olan $T(e,x,s)$ bağıntısı.
\item Kodu $s$ olan hesap dizisini bu hesabın son sonucuna götüren $U(s)$
  fonksiyonu.
\end{enumerate}
Aslında $T$ bağıntısı ile $U$ fonksiyonu ilkel özyinelemelidir.
\end{proof}

\begin{explain}
Normal Biçim Teoremi'nin titiz bir ispatını vermek için bir hesaplama modeli
saptamamız; hesaplanabilir fonksiyon betimlerinin ve hesap dizilerinin
kodlanmasını ayrıntılı biçimde gerçekleştirmemiz; ardından $T$ bağıntısı ile
$U$ fonksiyonunun ilkel özyinelemeli olduğunu doğrulamamız gerekirdi. Çoğu
uygulama için $T$ ile $U$'nun hesaplanabilir ve $U$'nun tümel olması yeter.

Normal biçim teoreminin ispatı belki de en iyi şu sloganla anımsanır:
$\umin{s}{T(e,x,s)}$, indisi $e$ olan fonksiyonun $x$ girdisindeki hesap
dizisini arar; $U$ ise böyle bir dizi bulunursa çıktısını döndürür.
\end{explain}

Hesaplama modeli kısmi özyinelemeli fonksiyonlarsa, Kleene'nin başlangıçta
kullandığı model de buydu, teorem herhangi bir fonksiyonun tanımı için
sınırsız aramanın yalnızca bir kez, yani tek bir
$\umin{y}{f(\vec x,y)}$ işleciyle kullanılmasının yeterli olduğunu gösterir.
Bu anlamda her kısmi özyinelemeli fonksiyonun bir normal biçimi vardır.

$T$ ile $U$, $\cfind{0},\cfind{1},\cfind{2},\dots$ numaralandırmasını
tanımlamak için kullanılabilir. Bundan sonra uygun bir $T$ ve $U$ seçimini
sabitlediğimizi varsayacak ve
\[
  \cfind{e}(x)\simeq U(\umin{s}{T(e,x,s)})
\]
denklemini $\cfind{e}$ fonksiyonunun \emph{tanımı} sayacağız.

Bir başka yararlı olgu şudur:

\begin{thm}
Her kısmi hesaplanabilir fonksiyonun sonsuz sayıda indisi vardır.
\end{thm}

Bu da sezgisel olarak açıktır. Hesaplanabilir bir fonksiyonun herhangi bir
betimi verildiğinde, aynı fonksiyonu, yani aynı girdi-çıktı çiftlerini
hesaplayan fakat önce hesaba etkisi olmayan bir işlem yapan, örneğin
$0=0$ mı diye sınayan ya da $5$'e kadar sayan farklı bir betim oluşturulur.
Değiştirilmiş betimin indisi her zaman ilk indisten farklıdır. İkisi aynı
fonksiyonun yalnızca biraz farklı hesaplanmış indisleridir.

